\documentclass[11pt,a4paper]{article}
\usepackage{fontspec}
\setmainfont{Noto Serif}
\setsansfont{Noto Sans}
\setmonofont{Noto Sans Mono}[Scale=0.86]
\usepackage[vietnamese]{babel}
\usepackage{amsmath,amssymb,amsthm,mathtools}
\usepackage{geometry,booktabs,longtable,tabularx,array}
\usepackage{tikz}
\usetikzlibrary{arrows.meta,positioning,calc}
\usepackage{xcolor,hyperref,xurl,fancyhdr,enumitem}
\geometry{left=2.25cm,right=2.25cm,top=2.3cm,bottom=2.3cm,headheight=15pt}
\definecolor{ink}{HTML}{183449}
\definecolor{teal}{HTML}{187C80}
\definecolor{pale}{HTML}{E8F2F2}
\hypersetup{colorlinks=true,linkcolor=ink,citecolor=teal,urlcolor=teal,
 pdftitle={Week 1: Project Brief và nền tảng nghiên cứu},
 pdfauthor={ECDR Research}}
\urlstyle{same}
\setlength{\parindent}{0pt}
\setlength{\parskip}{5pt}
\setlength{\emergencystretch}{3em}
\setlength{\tabcolsep}{4pt}
\renewcommand{\arraystretch}{1.15}
\setlist{nosep,leftmargin=1.5em}
\newcolumntype{Y}{>{\raggedright\arraybackslash}X}
\newcommand{\R}{\mathbb R}
\newcommand{\N}{\mathbb N}
\newcommand{\norm}[1]{\left\lVert #1\right\rVert}
\newcommand{\abs}[1]{\left\lvert #1\right\rvert}
\newcommand{\code}[1]{{\small\ttfamily\nolinkurl{#1}}}
\newcommand{\diag}{\operatorname{diag}}
\newcommand{\vecop}{\operatorname{vec}}
\newcommand{\relu}{\operatorname{ReLU}}
\newcommand{\ind}{\mathbf 1}
\newcommand{\E}{\mathbb E}
\newcommand{\Sset}{\mathcal S}
\newcommand{\Uset}{\mathcal U}
\newtheorem{definition}{Định nghĩa}[section]
\newtheorem{theorem}[definition]{Định lý}
\newtheorem{proposition}[definition]{Mệnh đề}
\newtheorem{lemma}[definition]{Bổ đề}
\newtheorem{corollary}[definition]{Hệ quả}
\theoremstyle{remark}
\newtheorem{example}[definition]{Ví dụ}
\newtheorem{remark}[definition]{Nhận xét}
\renewcommand{\proofname}{Chứng minh}
\pagestyle{fancy}
\fancyhf{}
\fancyhead[L]{\small\textcolor{ink}{ECDR | Cơ sở nghiên cứu và lý thuyết}}
\fancyhead[R]{\small Week 1}
\fancyfoot[L]{\small 28/10/2026}
\fancyfoot[R]{\thepage}
\begin{document}
\title{\textcolor{ink}{ECDR}\\[5pt]\Large Week 1: Project Brief và nền tảng nghiên cứu}
\author{Báo cáo nghiên cứu lý thuyết}
\date{28 tháng 10 năm 2026}
\maketitle
\noindent Bản tuần độc lập, trích từ báo cáo tổng hợp. Scope: lý thuyết mới
theo brief của người nghiên cứu; không dùng implementation cũ làm nguồn
định nghĩa phương pháp. Các tính toán minh họa không phải benchmark.
\tableofcontents
\clearpage
\part{Week 1 --- Project Brief và nền tảng nghiên cứu}
\section{Project brief}
\subsection{Bối cảnh và vấn đề nghiên cứu}
Một thiết bị nhận chuỗi cảm biến liên tục, cắt thành các cửa sổ $X_t$ rồi
chạy classifier $f_\theta$ cho từng cửa sổ. Hai cửa sổ có thể chứa nhiều mẫu
chung nhưng baseline vẫn đánh giá đầy đủ mô hình. Điều này đặt ra cơ hội
nghiên cứu: liệu một điều kiện kiểm tra đơn giản có thể xác định rằng quyết
định hiện tại đã được bao hàm bởi một chứng nhận trước đó?

Gọi $\mathcal P$ là preprocessing cố định và
\begin{equation}
 F=f_\theta\circ\mathcal P:\R^{C\times W}\longrightarrow\R^K,
 \qquad h(X)=\arg\max_{1\le k\le K}F_k(X).
 \label{eq:reference}
\end{equation}
Ở đây, $X\in\R^{C\times W}$ là cửa sổ tín hiệu thô có $C=6$ kênh (3 trục gia tốc kế $a_x,a_y,a_z$ đo bằng $g$ và 3 trục con quay hồi chuyển $g_x,g_y,g_z$ đo bằng $\text{rad/s}$) với độ rộng $W=128$ mẫu (tương đương $2{,}56\text{ s}$ ở tần số $50\text{ Hz}$). Ánh xạ $\mathcal P$ là phép tiền xử lý chuẩn hóa z-score cố định từng kênh $Z_{c,i}=\mathcal P(X)_{c,i}=(X_{c,i}-\mu_c)/\sigma_c$, với các tham số thống kê $\mu_c,\sigma_c$ được đóng băng hoàn toàn từ tập huấn luyện (training split). Tính chất \emph{cố định} (fixed) của $\mathcal P$ là điều kiện bắt buộc tiên quyết: nếu preprocessing bị thay đổi động giữa các cửa sổ (chẳng hạn như chuẩn hóa min-max cục bộ theo từng cửa sổ hay trung bình động streaming), ánh xạ tham chiếu $F$ sẽ bị dịch chuyển và chứng chỉ an toàn $\Sset_r$ xây dựng tại anchor sẽ bị vô hiệu hóa. Bộ phân loại $f_\theta$ là mạng nơ-ron sâu (mô hình SmallCNN1D) nhận tensor chuẩn hóa không thứ nguyên $Z=\mathcal P(X)\in\R^{C\times W}$ và xuất ra vector logits $K$ chiều ($K=6$ hoặc $K=12$). Toàn bộ chuỗi suy luận tham chiếu $F=f_\theta\circ\mathcal P$ ánh xạ trực tiếp từ tín hiệu vật lý thô $X$ sang logits, và $h(X)$ là quyết định lớp có logit lớn nhất. Quy tắc xử lý tie
phải cố định. Để đơn giản và bảo thủ, certificate trong báo cáo chỉ nhận
trường hợp lớp thắng duy nhất. Một \emph{trusted anchor} $X^r$ là cửa sổ đã
thực sự được đánh giá bởi reference, với $c=h(X^r)$.

\textbf{Problem statement.} Xây dựng một họ miền $\Sset_r$ và cơ chế kiểm tra
sao cho mọi query được chấp nhận thỏa
\begin{equation}
 X\in\Sset_r\quad\Longrightarrow\quad h(X)=h(X^r)=c,
 \label{eq:target}
\end{equation}
đồng thời xác định các điều kiện để chi phí xây dựng, kiểm tra và duy trì
certificate có thể được bù bởi các lần tránh chạy reference.

Mục tiêu đúng đắn là bảo toàn \emph{quyết định reference}. Nếu reference
phân loại sai ground truth, ECDR vẫn có thể bảo toàn quyết định sai đó. Trả
lại lớp $c$ cũng không đồng nghĩa tái tạo logits, xác suất softmax hoặc độ
tin cậy của query.

\subsection{Mục tiêu, đầu vào, đầu ra và ranh giới}
\begin{table}[!htbp]\centering
\caption{Project brief chuẩn hóa.}
\begin{tabularx}{\linewidth}{@{}p{2.5cm}Y@{}}\toprule
Thành phần & Đặc tả\\\midrule
O1: correctness & Chấp nhận reuse chỉ khi certificate hợp lệ và query thuộc
đúng miền đã chứng nhận.\\
O2: hiệu quả & Suy ra điều kiện hòa vốn; sau này kiểm bằng runtime đầy đủ.\\
O3: đối chiếu & So sánh guarantee, hình học miền và cost với các baseline;
LipCache là một đối chiếu cần kiểm đúng mô hình được chứng nhận.\\
Input & Chuỗi đa kênh có recording/subject/time identity, reference cố định,
anchor và candidate domain.\\
Output & Lớp reference hoặc fallback; certificate record có giả thiết và miền.\\
Phạm vi lý thuyết & Định nghĩa data, dịch cửa sổ, bounds, soundness, lifecycle,
membership và chi phí.\\
Giới hạn & Chưa chứng minh coverage hữu ích, tốc độ trên HAPT, superiority
so với LipCache hoặc độ chính xác của CNN 12 lớp.\\\bottomrule
\end{tabularx}
\end{table}

\subsection{Câu hỏi nghiên cứu}
\begin{enumerate}[label=RQ\arabic*.]
\item Dữ liệu cửa sổ có cấu trúc xác định nào có thể dùng để giảm bài toán
chứng nhận, và cấu trúc nào chỉ là giả thuyết thống kê?
\item Với một reference và một candidate domain, điều kiện đủ nào bảo toàn
quyết định trên toàn miền?
\item Miền căn chỉnh theo thời gian hoặc theo event có thể tăng xác suất
query được chứng nhận so với miền đồng nhất hay không?
\item Certificate có thể dùng nhiều bước mà không làm sai identity/state
hay không, và cần mang theo thông tin nào?
\item Những mức hit, overhead, rebuild và horizon nào cho lợi ích chi phí?
\end{enumerate}
RQ1--RQ2 có thành phần chứng minh được từ mô hình toán. RQ3 và lợi ích thực
tế của RQ5 cần thí nghiệm. RQ4 có thể giải quyết soundness bằng quy tắc đủ
bảo thủ, nhưng usefulness của quy tắc đó vẫn là câu hỏi thực nghiệm.

\section{Thuật ngữ và cấu trúc của một chứng nhận}
\begin{definition}[Window, event, region và certificate]
Window là một ma trận quan sát $X\in\R^{C\times W}$. Event descriptor $e$
là thông tin dùng chọn một họ miền hoặc budget. Region $\Sset_r$ là tập
các đầu vào khả dĩ, không phải một đoạn thời gian. Certificate là bằng
chứng cho thuộc tính của $F$ trên toàn tập đó, kèm model/preprocessing/domain
identity và các giả thiết số học.
\end{definition}
Ví dụ, ``biến thiên thấp'' có thể chọn box hẹp hơn. Tên event đúng không
thay thế chứng nhận. Event được chọn sai vẫn an toàn nếu miền thực sự đã
được chứng nhận và query thật sự thuộc miền; tác động khi đó chủ yếu lên
coverage hoặc chi phí.

Ba tầng cần đi đúng thứ tự là
\begin{equation}
\underbrace{\text{candidate geometry}}_{\text{đề xuất tập}}
\longrightarrow
\underbrace{\text{whole-domain bound}}_{\text{chứng minh thuộc tính}}
\longrightarrow
\underbrace{\text{query membership}}_{\text{áp dụng cho điểm hiện tại}}.
\end{equation}
Một box có thể kiểm membership rất nhanh nhưng không giữ quyết định.
Một certificate đúng có thể chỉ bao phủ anchor. Một miền có nhiều query
có thể yêu cầu bound rất đắt. Do đó ba tầng không thể thay thế nhau.

\section{Literature và định vị giả thuyết đóng góp}
\subsection{Các nền tảng được kế thừa}
Lipschitz--margin liên hệ biên độ biến thiên đầu ra với khoảng cách đến
biên quyết định \cite{lmt}. IBP lan truyền interval qua các phép toán để
bao đầu ra \cite{ibp}. CROWN xây dựng các relaxation để thu được cận hợp lệ
\cite{crown}; auto-LiRPA mở rộng cách tính cận trên computational graph
\cite{lirpa}. Các cơ chế này là nền tảng kế thừa, không phải định lý mới của
ECDR. Phần chứng minh trong Week 3 trình bày các phiên bản đủ dùng và chỉ rõ
điều kiện áp dụng.

StreamiNNC nghiên cứu tái sử dụng activation khi CNN xử lý cửa sổ chồng
lấn, đồng thời phân tích khó khăn của padding và pooling \cite{streaminnc}.
TinyD\'ej\`aVu nghiên cứu tính tăng dần và bộ nhớ trên sensor streams/MCU
\cite{tiny}. Hai công trình làm cho claim ``đầu tiên khai thác overlap của
CNN'' không có cơ sở. Quy tắc exact cho CNN valid trong báo cáo là một
chuyên biệt hóa tự suy ra, không phải tuyên bố phát minh reuse activation.

LipCache là đối chiếu về certified caching \cite{lipcache}. Theo phạm vi
của paper, certificate nằm ở phía GuardNet; bảo toàn MainNet là một nghĩa
khác và không được tự suy ra. Vì thế bảng so sánh phải ghi \emph{model nào}
được bảo đảm và \emph{payload nào} được trả. Không thể chuyển một con số
speedup trên ảnh của paper thành tốc độ trên HAPT.

\begin{table}[!htbp]\centering\small
\caption{So sánh theo đối tượng và bảo đảm; chưa phải bảng benchmark.}
\begin{tabularx}{\linewidth}{@{}p{2.25cm}YYY@{}}\toprule
Hướng & Đối tượng chính & Điều được cung cấp & Vai trò với ECDR\\\midrule
Full CNN & Mỗi window & Quyết định reference sau forward & Baseline kết quả/cost\\
Heuristic reuse & Similarity/confidence & Quy tắc thực nghiệm; tự nó không là
whole-domain proof & Negative control\\
StreamiNNC & Streaming activations & Cơ chế giảm tính trùng; cần xử lý biên
và alignment & Prior art overlap\\
TinyD\'ej\`aVu & Incremental sensor inference & Cơ chế streaming và bộ nhớ
trong thiết lập paper & Prior art systems\\
LipCache & GuardNet/cache & Certificate theo phạm vi GuardNet paper & Đối
chiếu cache semantics\\
Lipschitz / IBP / CROWN & Mạng và miền perturbation & Cận có điều kiện trên miền
được khai báo & Machinery cho proof\\
ECDR đề xuất & Quyết định frozen reference & Reuse khi bound và membership
đúng; usefulness còn mở & Câu hỏi nghiên cứu\\\bottomrule
\end{tabularx}
\end{table}

\subsection{Research gap được phát biểu có thể kiểm tra}
Giả thuyết đóng góp là kết hợp cấu trúc cửa sổ với chứng nhận quyết định
reference và một policy có cost đủ thấp. Để claim này đứng được, cần chứng
minh đồng thời: miền được định nghĩa đúng với data; bound sound trên miền;
runtime membership đúng; lifecycle không làm drift certificate; và tổng
chi phí thấp hơn baseline trong một operating regime được công bố.
Hiện chưa có bằng chứng cho claim ``ECDR mới hoàn toàn'', ``luôn nhanh hơn''
hoặc ``tốt hơn LipCache''. Bảng trên là targeted related work, không được
gọi là systematic review đầy đủ.

\section{Kế hoạch nghiên cứu theo ba tuần}
\begin{table}[!htbp]\centering
\caption{Câu hỏi, đầu ra và tiêu chí hoàn thành về nội dung.}
\begin{tabularx}{\linewidth}{@{}p{1.35cm}YYY@{}}\toprule
Tuần & Câu hỏi & Đầu ra & Điều kiện hoàn thành\\\midrule
Week 1 & Nghiên cứu cái gì và bảo đảm gì? & Brief, RQ, thuật ngữ, prior work,
gap, constraints & Phân biệt decision preservation, task accuracy và cost\\
Week 2 & Input thật sự là gì, liên quan thế nào? & Data contract, window/label,
normalization, shift, complexity & Truy được mỗi biến toán về đối tượng dữ liệu\\
Week 3 & Khi nào reuse đúng và hữu ích? & Contract, proofs, counterexamples,
miền có cấu trúc, amortization, hypotheses & Mỗi claim có giả thiết, lập luận
và phần chưa chứng minh\\\bottomrule
\end{tabularx}
\end{table}
Đây là cấu trúc trình bày nghiên cứu theo bản chuẩn hóa mới của người nghiên
cứu, không phải lịch thực nghiệm được phục dựng từ timestamp. Hoàn thành
báo cáo lý thuyết không tự hoàn thành verifier hoặc đo runtime.

\textbf{Kết quả Week 1.} Bài toán đã được đóng thành một thuộc tính kiểm
chứng được trên reference cố định, một mục tiêu chi phí riêng và năm RQ.
Week 2 phải xác định cấu trúc dữ liệu đủ chặt để Week 3 không xây proof trên
một biểu diễn hoặc preprocessing không tồn tại.

\section*{Nguồn và cách tái kiểm}
Nội dung theo bản chuẩn hóa Week 1--3 do người nghiên cứu cung cấp và các
nguồn gốc trong bibliography. Các số data mô tả snapshot local theo policy
đã ghi; mô hình lý thuyết không phải checkpoint triển khai cũ. Báo cáo tổng
hợp trong cùng thư mục có đầy đủ bảng ký hiệu, provenance và các tuần còn
lại. Phép tính số học được lưu trong \code{recompute.py} và
\code{calculations.json}. Không có training hoặc benchmark ECDR mới.

\begingroup\small
\begin{thebibliography}{99}
\bibitem{hapt} J. Reyes-Ortiz, D. Anguita, L. Oneto, X. Parra.
\emph{Smartphone-Based Recognition of Human Activities and Postural
Transitions}. UCI Machine Learning Repository, 2015.
DOI: \href{https://doi.org/10.24432/C54G7M}{10.24432/C54G7M}.
Metadata và README local được đối chiếu ngày 28/10/2026.

\bibitem{lmt} Y. Tsuzuku, I. Sato, M. Sugiyama.
\emph{Lipschitz-Margin Training: Scalable Certification of Perturbation
Invariance for Deep Neural Networks}. NeurIPS, 2018.
\url{https://arxiv.org/abs/1802.04034}.

\bibitem{ibp} S. Gowal et al.
\emph{On the Effectiveness of Interval Bound Propagation for Training
Verifiably Robust Models}. arXiv:1810.12715, 2018.
\url{https://arxiv.org/abs/1810.12715}.

\bibitem{crown} H. Zhang, T.-W. Weng, P.-Y. Chen, C.-J. Hsieh, L. Daniel.
\emph{Efficient Neural Network Robustness Certification with General
Activation Functions}. NeurIPS, 2018.
\url{https://arxiv.org/abs/1811.00866}.

\bibitem{lirpa} K. Xu et al.
\emph{Automatic Perturbation Analysis for Scalable Certified Robustness
and Beyond}. NeurIPS, 2020.
\url{https://arxiv.org/abs/2002.12920}.

\bibitem{crownibp} H. Zhang et al.
\emph{Towards Stable and Efficient Training of Verifiably Robust Neural
Networks}. NeurIPS, 2019.
\url{https://arxiv.org/abs/1906.06316}.

\bibitem{betacrown} S. Wang et al.
\emph{Beta-CROWN: Efficient Bound Propagation with Per-neuron Split
Constraints for Neural Network Robustness Verification}. NeurIPS, 2021.
\url{https://papers.neurips.cc/paper_files/paper/2021/hash/fac7fead96dafceaf80c1daffeae82a4-Abstract.html}.

\bibitem{montufar} G. Mont\'ufar, R. Pascanu, K. Cho, Y. Bengio.
\emph{On the Number of Linear Regions of Deep Neural Networks}.
NeurIPS, 2014.
\url{https://proceedings.neurips.cc/paper_files/paper/2014/hash/fa6f2a469cc4d61a92d96e74617c3d2a-Abstract.html}.

\bibitem{streaminnc} C. Kechris, J. Dan, J. Miranda, D. Atienza.
\emph{Don't Think It Twice: Exploit Shift Invariance for Efficient Online
Streaming Inference of CNNs}. arXiv:2408.03223, 2024.
\url{https://arxiv.org/abs/2408.03223}.

\bibitem{tiny} Z. Huang, E. Baccelli.
\emph{TinyD\'ej\`aVu: Smaller Memory Footprint \& Faster Inference on Sensor
Data Streams with Always-On Microcontrollers}. arXiv:2512.09786, 2025.
\url{https://arxiv.org/abs/2512.09786}. Tên theo metadata arXiv được tra
28/10/2026; không nhập số benchmark của paper vào kết quả HAPT.

\bibitem{lipcache} Z. Xiang, Y. Chen, F. Ying, B. Zhou, H. Zhao, S. Dustdar.
\emph{LipCache: A Local Inference Proxy with Certified Caching for Edge
Image Classification Service}. arXiv:2608.13144, 2026.
\url{https://arxiv.org/abs/2608.13144}. Đối chiếu PDF v1 lưu local,
đặc biệt phạm vi GuardNet/MainNet tại phần certification và system semantics.
\end{thebibliography}

\endgroup

\end{document}
